\documentclass[10pt,a4paper]{amsart}
\usepackage[margin=19mm]{geometry}
\usepackage{amsmath,amssymb,mathtools}
\usepackage[scr=rsfs,cal=euler]{mathalfa}
\usepackage{xfrac}
\usepackage[unicode,hidelinks,bookmarks=false]{hyperref}
\newcommand{\ie}{i.e.\ }
\newcommand{\cf}{cf.\ }
\newcommand{\resp}{respectively\ }
\newcommand{\Subsection}{\subsection}
\providecommand{\nobreakdash}{\nobreak\hbox{-}}
\setlength{\emergencystretch}{2em}
\multlinegap0pt
\usepackage{bm}
\usepackage{bbm}
\usepackage{dsfont}
\usepackage{xspace}
\usepackage{cancel}
\usepackage{mathtools}
\usepackage{amsthm}
\makeatletter
\@ifundefined{theorem}{\newtheorem{theorem}{Theorem}}{}
\@ifundefined{lemma}{\newtheorem{lemma}[theorem]{Lemma}}{}
\makeatother
\title[The Condition for Structured Coding to Improve Random Coding]{The Condition for Structured Coding to Improve Random Coding in the Binary Modulo-sum Problem}
\author{Yohsuke Tsujino, Shun Watanabe}
\date{Source: arXiv:2606.25570v2 (CC BY 4.0)}
\begin{document}
\maketitle

\noindent\emph{Source and licence.} The Condition for Structured Coding to Improve Random Coding in the Binary Modulo-sum Problem. Version arXiv:2606.25570v2 (CC BY 4.0), distributed under the Creative Commons Attribution 4.0 International licence, as recorded in the arXiv licence metadata for this exact version and shown on the abstract page https://arxiv.org/abs/2606.25570v2. Full rights notice: licenses/arxiv-2606.25570-CC-BY-4.0.txt.\par\smallskip

\noindent\textit{Editorial scope.} Counted unit: the Lemma whose source label is \texttt{lem:boundary} in the article (source0.tex lines 910--920, source label \texttt{lem:boundary}), delivered together with the article's own complete original proof of it (source0.tex lines 922--986, one \texttt{proof} environment of 65 lines). No mathematics is written, repaired or re-derived: statement and proof are the article's own text line for line. Required context retained uncounted: none. Every definition, display and label of those lines is retained unchanged. Omitted from this package and not redistributed: 1--909, 921--921, 987--1558. Nothing inside a retained excerpt is omitted or altered: every retained body is delivered exactly as the article's lines are written, so no omission receipt of any kind accompanies this package. Citations: the retained excerpts carry no citation command at all, so no bibliography entry of the article is transcribed and no reference is redistributed. Reference adaptation: none was needed and none was made; every reference inside the delivered unit resolves inside it, so no unresolved reference or citation is produced. Printed slips kept literal: no printed word, symbol, hypothesis or number of the counted statement or of its proof was altered. Layout and class: the article's own class is \texttt{IEEEtran} with packages including amsmath, amssymb, amsthm, enumitem, standalone, bm, graphics, mathtools; the wrapper is the lane's pinned \texttt{amsart} editorial wrapper at 10pt on A4, which restates the theorem-like environments the retained text needs and adds the article's own macro definitions that the retained text needs (none), with extra wrapper packages (bm, bbm, dsfont, xspace, cancel, mathtools). The delivered numbering is the unit's own. Assets: no retained span contains a figure environment, a graphics-inclusion command, a PDF-page inclusion, a table or a TikZ picture; no image file is copied into this package, the package declares no asset dependency and no image file is redistributed with it. Licence: version arXiv:2606.25570v2 of this article is distributed under the Creative Commons Attribution 4.0 International licence, as recorded in the arXiv licence metadata for this exact version and shown on the abstract page https://arxiv.org/abs/2606.25570v2. A full rights notice is delivered at licenses/arxiv-2606.25570-CC-BY-4.0.txt. Characters: every retained excerpt is pure ASCII, so the delivered engine, XeLaTeX with its default Latin Modern text font, reports no missing character.
    \begin{lemma}\label{lem:boundary}
    Let $P$ and $P'$ be two probability distributions on the same finite alphabet and suppose that $\operatorname{supp}(P) = \operatorname{supp}(P')$.
    Consider the minimization of $D(Q\|P')$ over all probability distributions $Q$ satisfying
    \begin{align*}
        D(Q\|P)\leq D(Q\|P').
    \end{align*}
    Then the minimum is attained on the boundary
    \begin{align*}
        D(Q\|P)=D(Q\|P').
    \end{align*}
    \end{lemma}

    \begin{proof}
    The case $P=P'$ is trivial, so we assume that $P\neq P'$.
    Suppose, for contradiction, that the minimum is attained at a distribution $Q^*$ satisfying the strict inequality
    \begin{align*}
        D(Q^*\|P)<D(Q^*\|P').
    \end{align*}
    
    For $0 \leq \lambda \leq 1$, define
    \begin{align*}
        Q_\lambda=\lambda Q^*+(1-\lambda)P'.
    \end{align*}
    By the convexity of divergence in the first argument, for $\lambda \in (0,1)$, we have
    \begin{align}
        D(Q_\lambda\|P')&\leq \lambda D(Q^*\|P')+(1-\lambda)D(P'\|P')\notag\\ 
        &=\lambda D(Q^*\|P')\notag\\
        &< D(Q^*\|P'). \label{eq:divergence to P'}
    \end{align}
    % Thus, every distribution on the line segment between $Q^*$ and $Q^{**}$ has divergence to $P'$ no larger than that of $Q^*$.

    Now define
    \begin{align*}
        F(\lambda)=D(Q_\lambda\|P)-D(Q_\lambda\|P').
    \end{align*}
    Since \(D(Q^*\|P')<\infty\), we have
    \[
    \operatorname{supp}(Q^*)\subseteq \operatorname{supp}(P').
    \]
    By the assumption \(\operatorname{supp}(P)=\operatorname{supp}(P')\),
    it follows that
    \[
    \operatorname{supp}(Q_\lambda)\subseteq \operatorname{supp}(P)
    =\operatorname{supp}(P')
    \]
    for every \(\lambda\in[0,1]\). Therefore, both
    \(D(Q_\lambda\|P)\) and \(D(Q_\lambda\|P')\) are finite and continuous
    in \(\lambda\). Hence, \(F(\lambda)\) is continuous on \([0,1]\).
    Since $\operatorname{supp}(P)=\operatorname{supp}(P')$, we have $D(P'\|P)<\infty$.
    Moreover,
    \begin{align*}
        F(1)=D(Q^*\|P)-D(Q^*\|P')<0,
    \end{align*}
    whereas
    \begin{align*}
        F(0)=D(P'\|P)-D(P'\|P')=D(P'\|P)>0.
    \end{align*}
    Hence, by the intermediate value theorem, there exists $\lambda_0\in(0,1)$ such that
    \begin{align*}
        F(\lambda_0)=0.
    \end{align*}
    That is,
    \begin{align*}
        D(Q_{\lambda_0}\|P)=D(Q_{\lambda_0}\|P').
    \end{align*}
    Furthermore,
    \begin{align*}
        D(Q_{\lambda_0}\|P') < D(Q^*\|P')
    \end{align*}
    by \eqref{eq:divergence to P'}.
    Therefore, there exists a feasible distribution on the boundary whose objective value is strictly smaller than that of $Q^*$.
    This contradicts that $Q^*$ is the minimizer. 
    Thus, the minimum is attained on the boundary
    \begin{align*}
        D(Q\|P)=D(Q\|P').
    \end{align*}
    \end{proof}

\end{document}
